\chapter{RESULTS AND FINDINGS}

\section{Interpretation of Lexical vs. Semantic Trade-Offs}
However, there is an undeniably explicit story behind the aggregated data obtained in the benchmarking stage. The implementation, that is Approach A, is capable of performing fast deterministic resume filtering. But, Benchmark 4 explicitly highlights the ``semantic gap.'' Given the use of Regular Expressions and term-frequency statistics in this approach, the candidate and the recruiter have to use precisely the same terminology. In case of any discrepancy in the usage of ``REST'' vs ``RESTful'' or ``Postgres'' vs ``PostgreSQL,'' then no points are given for that particular keyword within these constraints.

On the other hand, Sentence-BERT and LLMs (Approach C) intrinsically know the mappings between these terms, meaning they fall into the identical vector space. Deshmukh \& Raut explain this concept in the relevant literature. But, the justification for avoiding using such architecture within our implementation (Approach C) is given by the results obtained from Benchmark 5.

\section{Viability of the TF-IDF and Regex Approach}
True effectiveness of the prototype is determined solely by how the system will perform in the desired setting. Within an actual SaaS business environment, for instance, when a junior software developer job opening generates 4,000 applications in a day, scanning them using an LLM at a rate of 4.18 seconds per document would require 4.5 hours of uninterrupted heavy-GPU computation. On the other hand, utilizing the prototype, all 4,000 resumes would be effectively parsed, scored, and ranked within just 25 seconds using a typical low-budget CPU.

Clearly, TF-IDF and regex are both very practical and relevant solutions in their role as ``first-pass filters'' that enable the discarding of the bottom 80\% of applicants whose resumes have no corresponding technical qualifications within a relatively short period of time. The HR team will then be able to transfer the remaining candidates to either semantic engines or to human assessment.

\section{Insights into Transparency and Auditability}
One important insight is the successful operation of the rule-based Explainable AI module, verified in Benchmark 3. Despite the high sophistication level of such explainable algorithms as SHAP, their resulting graphs and perturbation tables can mislead non-technical HR managers, which is contrary to the basic concept of an "explanation." The prototype methodology of capturing the mathematical score dictionary and providing an unambiguous text paragraph is a guaranteed solution. Should the candidate ask about the reasoning behind the rejection, the HR manager would instantly answer him using the provided text: "You met 50\% of required skills, but you are missing the PostgreSQL skill set." Such an approach will protect the company from the potential legal bias issues and instill confidence in the software layer within the company itself.

\section{Implications for Practical Recruitment}
It is evident that the prototype demonstrates how algorithmic fairness and speed are not necessarily contradictory concepts. Standardization of the process via using Tesseract OCR helps prevent visual format biases, meaning that applicants from any social status background without access to paid resume-formating services will be assessed based solely on the textual content.
